\documentclass[11pt,a4paper]{article}
\newcommand{\intestazione}{Dimero di spin --- il filo conduttore}
% =====================================================================
%  Preambolo comune ai documenti sul dimero (serie "che cosa abbiamo fatto")
%  Incluso con % =====================================================================
%  Preambolo comune ai documenti sul dimero (serie "che cosa abbiamo fatto")
%  Incluso con % =====================================================================
%  Preambolo comune ai documenti sul dimero (serie "che cosa abbiamo fatto")
%  Incluso con \input{preambolo_dimero} subito dopo \documentclass.
% =====================================================================
\usepackage[T1]{fontenc}
\usepackage[utf8]{inputenc}
\usepackage[main=italian,provide=*]{babel}
\usepackage{amsmath,amssymb,mathtools}
\usepackage{graphicx}
\usepackage{booktabs}
\usepackage{colortbl}
\usepackage{xcolor}
\usepackage{geometry}
\usepackage{placeins}
\usepackage{caption}
\usepackage{enumitem}
\usepackage{fancyhdr}
\usepackage{needspace}
\usepackage{tikz}
\usetikzlibrary{arrows.meta,positioning,shapes.geometric,calc,fit,backgrounds}
\usepackage{hyperref}
\hypersetup{colorlinks=true,linkcolor=blu,citecolor=blu,urlcolor=blu}

\geometry{margin=2.5cm}

% --- colori coerenti con quelli delle figure (palette Okabe-Ito)
\definecolor{blu}{HTML}{0072B2}
\definecolor{arancio}{HTML}{E69F00}
\definecolor{verdes}{HTML}{009E73}
\definecolor{rossos}{HTML}{D55E00}
\definecolor{violas}{HTML}{CC79A7}
\definecolor{grigetto}{HTML}{E8E8E8}

\captionsetup{font=small,labelfont=bf,width=0.92\textwidth}

\pagestyle{fancy}
\fancyhf{}
\fancyhead[L]{\small\itshape\intestazione}
\fancyhead[R]{\small\thepage}
\renewcommand{\headrulewidth}{0.3pt}

% --- notazione
\newcommand{\ket}[1]{\lvert #1\rangle}
\newcommand{\bra}[1]{\langle #1 \rvert}
\newcommand{\media}[1]{\langle #1 \rangle}
\newcommand{\vs}[1]{\mathbf{s}_{#1}}

% --- riquadro "in una riga": il messaggio da portare a casa
\newcommand{\messaggio}[1]{%
  \begin{center}
  \begin{minipage}{0.93\textwidth}
  \setlength{\fboxsep}{9pt}%
  \colorbox{grigetto}{\begin{minipage}{\dimexpr\textwidth-2\fboxsep}
  \small\textbf{In una riga.}\enspace #1
  \end{minipage}}
  \end{minipage}
  \end{center}
}

% --- riquadro "come lo sappiamo": la verifica numerica dietro un'affermazione
\newcommand{\verifica}[1]{%
  \begin{center}
  \begin{minipage}{0.93\textwidth}
  \small\textcolor{verdes}{\rule{2.2em}{1.1pt}}\enspace
  \textcolor{verdes}{\textbf{Come lo sappiamo.}}\enspace #1
  \end{minipage}
  \end{center}
  \medskip
}
 subito dopo \documentclass.
% =====================================================================
\usepackage[T1]{fontenc}
\usepackage[utf8]{inputenc}
\usepackage[main=italian,provide=*]{babel}
\usepackage{amsmath,amssymb,mathtools}
\usepackage{graphicx}
\usepackage{booktabs}
\usepackage{colortbl}
\usepackage{xcolor}
\usepackage{geometry}
\usepackage{placeins}
\usepackage{caption}
\usepackage{enumitem}
\usepackage{fancyhdr}
\usepackage{needspace}
\usepackage{tikz}
\usetikzlibrary{arrows.meta,positioning,shapes.geometric,calc,fit,backgrounds}
\usepackage{hyperref}
\hypersetup{colorlinks=true,linkcolor=blu,citecolor=blu,urlcolor=blu}

\geometry{margin=2.5cm}

% --- colori coerenti con quelli delle figure (palette Okabe-Ito)
\definecolor{blu}{HTML}{0072B2}
\definecolor{arancio}{HTML}{E69F00}
\definecolor{verdes}{HTML}{009E73}
\definecolor{rossos}{HTML}{D55E00}
\definecolor{violas}{HTML}{CC79A7}
\definecolor{grigetto}{HTML}{E8E8E8}

\captionsetup{font=small,labelfont=bf,width=0.92\textwidth}

\pagestyle{fancy}
\fancyhf{}
\fancyhead[L]{\small\itshape\intestazione}
\fancyhead[R]{\small\thepage}
\renewcommand{\headrulewidth}{0.3pt}

% --- notazione
\newcommand{\ket}[1]{\lvert #1\rangle}
\newcommand{\bra}[1]{\langle #1 \rvert}
\newcommand{\media}[1]{\langle #1 \rangle}
\newcommand{\vs}[1]{\mathbf{s}_{#1}}

% --- riquadro "in una riga": il messaggio da portare a casa
\newcommand{\messaggio}[1]{%
  \begin{center}
  \begin{minipage}{0.93\textwidth}
  \setlength{\fboxsep}{9pt}%
  \colorbox{grigetto}{\begin{minipage}{\dimexpr\textwidth-2\fboxsep}
  \small\textbf{In una riga.}\enspace #1
  \end{minipage}}
  \end{minipage}
  \end{center}
}

% --- riquadro "come lo sappiamo": la verifica numerica dietro un'affermazione
\newcommand{\verifica}[1]{%
  \begin{center}
  \begin{minipage}{0.93\textwidth}
  \small\textcolor{verdes}{\rule{2.2em}{1.1pt}}\enspace
  \textcolor{verdes}{\textbf{Come lo sappiamo.}}\enspace #1
  \end{minipage}
  \end{center}
  \medskip
}
 subito dopo \documentclass.
% =====================================================================
\usepackage[T1]{fontenc}
\usepackage[utf8]{inputenc}
\usepackage[main=italian,provide=*]{babel}
\usepackage{amsmath,amssymb,mathtools}
\usepackage{graphicx}
\usepackage{booktabs}
\usepackage{colortbl}
\usepackage{xcolor}
\usepackage{geometry}
\usepackage{placeins}
\usepackage{caption}
\usepackage{enumitem}
\usepackage{fancyhdr}
\usepackage{needspace}
\usepackage{tikz}
\usetikzlibrary{arrows.meta,positioning,shapes.geometric,calc,fit,backgrounds}
\usepackage{hyperref}
\hypersetup{colorlinks=true,linkcolor=blu,citecolor=blu,urlcolor=blu}

\geometry{margin=2.5cm}

% --- colori coerenti con quelli delle figure (palette Okabe-Ito)
\definecolor{blu}{HTML}{0072B2}
\definecolor{arancio}{HTML}{E69F00}
\definecolor{verdes}{HTML}{009E73}
\definecolor{rossos}{HTML}{D55E00}
\definecolor{violas}{HTML}{CC79A7}
\definecolor{grigetto}{HTML}{E8E8E8}

\captionsetup{font=small,labelfont=bf,width=0.92\textwidth}

\pagestyle{fancy}
\fancyhf{}
\fancyhead[L]{\small\itshape\intestazione}
\fancyhead[R]{\small\thepage}
\renewcommand{\headrulewidth}{0.3pt}

% --- notazione
\newcommand{\ket}[1]{\lvert #1\rangle}
\newcommand{\bra}[1]{\langle #1 \rvert}
\newcommand{\media}[1]{\langle #1 \rangle}
\newcommand{\vs}[1]{\mathbf{s}_{#1}}

% --- riquadro "in una riga": il messaggio da portare a casa
\newcommand{\messaggio}[1]{%
  \begin{center}
  \begin{minipage}{0.93\textwidth}
  \setlength{\fboxsep}{9pt}%
  \colorbox{grigetto}{\begin{minipage}{\dimexpr\textwidth-2\fboxsep}
  \small\textbf{In una riga.}\enspace #1
  \end{minipage}}
  \end{minipage}
  \end{center}
}

% --- riquadro "come lo sappiamo": la verifica numerica dietro un'affermazione
\newcommand{\verifica}[1]{%
  \begin{center}
  \begin{minipage}{0.93\textwidth}
  \small\textcolor{verdes}{\rule{2.2em}{1.1pt}}\enspace
  \textcolor{verdes}{\textbf{Come lo sappiamo.}}\enspace #1
  \end{minipage}
  \end{center}
  \medskip
}


\title{\bfseries Il dimero di spin: che cosa abbiamo fatto, e perché\\[2pt]
\large Documento 0 --- il filo conduttore}
\author{Samuele Schiavi \\ \small Tesi triennale in Fisica, Università di Parma
--- Relatore: Prof.\ Alessandro Chiesa}
\date{}

\begin{document}
\maketitle
\thispagestyle{fancy}

\begin{abstract}
\noindent
Questa serie di documenti racconta il lavoro svolto sul sistema più piccolo del
progetto --- due spin $1/2$ accoppiati --- seguendo l'ordine in cui è stato
effettivamente fatto, non l'ordine in cui lo si dedurrebbe dalla teoria. Ogni
affermazione numerica riportata è stata ricalcolata eseguendo il codice del
progetto, non ripresa da documenti precedenti. La teoria compare solo dove serve
a capire perché un passo viene dopo l'altro: nessuna derivazione, solo il
risultato che motiva la scelta successiva. Il dimero non è un caso di studio a
sé: è il banco di prova su cui è stato messo a punto ogni strumento poi
riapplicato a $N=3$.
\end{abstract}

\section{Perché due spin}

Due spin $1/2$ danno uno spazio di Hilbert di dimensione $4$. È piccolo abbastanza
da poter essere diagonalizzato esattamente --- a mano, in forma chiusa --- e
grande abbastanza da contenere tutti gli ingredienti fisici che contano:
entanglement, un incrocio di livelli al variare del campo, e la possibilità di
rompere una simmetria per anticiparlo.

Questo è precisamente ciò che serve per validare un metodo. Ogni volta che il
calcolo quantistico produce un numero, esiste un numero esatto con cui
confrontarlo. Quando poi lo stesso codice viene esteso a $N=3$, dove la forma
chiusa diventa più laboriosa, si sa già che gli strumenti funzionano perché sono
stati verificati dove la risposta era nota.

\messaggio{Il dimero non serve a scoprire fisica nuova: serve a stabilire che il
metodo dà la risposta giusta là dove la risposta giusta la conosciamo già.}

\section{I quattro passi, e come si incastrano}

Il lavoro si articola in quattro blocchi, ciascuno con un proprio documento. Non
sono argomenti indipendenti messi in fila: ognuno consuma l'output del
precedente.

\begin{figure}[htbp]
\centering
\resizebox{\textwidth}{!}{%
\begin{tikzpicture}[
  scatola/.style={draw,rounded corners=3pt,minimum height=1.45cm,
                  minimum width=3.0cm,align=center,font=\small,thick},
  freccia/.style={-{Stealth[length=3mm]},thick,gray}]

\node[scatola,draw=blu,fill=blu!8]      (A) at (0,0)    {\textbf{1. Sistema}\\[1pt]
   spettro esatto\\ ruolo del termine DM};
\node[scatola,draw=verdes,fill=verdes!8](B) at (5.0,0)  {\textbf{2. VQE}\\[1pt]
   stato fondamentale\\ scelta dell'ansatz};
\node[scatola,draw=arancio,fill=arancio!10](C) at (10.0,0){\textbf{3. Dinamica}\\[1pt]
   $e^{-iHt}$ via Trotter\\ errore vs.\ passi};
\node[scatola,draw=rossos,fill=rossos!8](D) at (15.0,0) {\textbf{4. Correlatori}\\[1pt]
   ancilla, 36 combinazioni\\ contenuto spettrale};

\draw[freccia] (A) -- node[above,font=\scriptsize,text=black]{riferimento}
                        node[below,font=\scriptsize,text=black]{esatto} (B);
\draw[freccia] (B) -- node[above,font=\scriptsize,text=black]{stato}
                        node[below,font=\scriptsize,text=black]{iniziale} (C);
\draw[freccia] (C) -- node[above,font=\scriptsize,text=black]{$U(t)$}
                        node[below,font=\scriptsize,text=black]{nel circuito} (D);
\end{tikzpicture}}
\caption{La catena di dipendenze. Lo spettro esatto del passo 1 è il metro con
cui si giudica il VQE; lo stato fondamentale del passo 2 è lo stato di partenza
dell'evoluzione temporale; l'evoluzione del passo 3 è un blocco interno al
circuito dei correlatori. Un errore in un passo si propaga silenziosamente a
tutti i successivi: è la ragione per cui ogni blocco è stato validato
separatamente prima di passare al seguente.}
\end{figure}

\FloatBarrier

\section{L'Hamiltoniana, e il solo pezzo di teoria che serve ricordare}

Il modello è un dimero di Heisenberg isotropo in campo, con l'aggiunta di un
termine antisimmetrico di Dzyaloshinskii--Moriya (DM):
\begin{equation}
H = \underbrace{J\,(X_1X_2+Y_1Y_2+Z_1Z_2)}_{\text{scambio isotropo}}
  + \underbrace{b\,(Z_1+Z_2)}_{\text{campo}}
  + \underbrace{D\,(X_1Z_2-Z_1X_2)}_{\text{termine DM}}.
\label{eq:H}
\end{equation}
Ogni spin $1/2$ si mappa direttamente su un qubit: nessuna trasformazione
fermionica, nessun mapping intermedio. Questo è il motivo per cui i sistemi
molecolari di spin sono un banco di prova naturale per l'hardware quantistico
attuale --- il costo di rappresentazione è minimo.

L'unico fatto teorico che serve tenere a mente in tutta la serie è il seguente.
Con $D=0$, l'Hamiltoniana commuta con la magnetizzazione totale
$M=(Z_1+Z_2)/2$: i suoi autostati hanno un valore definito di $M$, e stati con
$M$ diverso non possono mescolarsi. Con $D\neq0$ questo cade. Tutto ciò che
segue --- la forma dello spettro, il fallimento di un ansatz, la struttura dei
correlatori --- discende da questa singola frase.

\messaggio{Il termine DM non aggiunge complicazione per il gusto di aggiungerla:
rompe una simmetria, e la rottura è il filo che collega i quattro documenti.}

\section{I due punti di lavoro usati}

Il sistema ha tre parametri, $(J,b,D)$; fissando $J=1$ come unità restano
$b/J$ e $D/J$. Nel corso del lavoro sono stati usati due punti distinti, per
ragioni diverse. Sono dichiarati qui una volta per tutte, così che nelle figure
dei documenti successivi non ci sia ambiguità.

\begin{table}[htbp]
\centering
\footnotesize
\renewcommand{\arraystretch}{1.35}
\begin{tabular}{@{}p{2.9cm}p{3.4cm}p{2.4cm}p{5.6cm}@{}}
\toprule
\rowcolor{grigetto}
\textbf{Nome} & \textbf{Parametri} & \textbf{Dove è usato} & \textbf{Perché quello} \\
\midrule
Scansione in campo & $D/J\in\{0,\,0.2\}$, $b/J$ variabile
  & Documenti 1 e 2 & Serve vedere l'incrocio, quindi $b$ deve variare \\
Punto ``test 2'' & $b/J=0.35$, $D/J=0.80$
  & Documento 4 & Stato fondamentale non banale, termine DM ben acceso \\
Punto dinamico & $b/J=-0.18$, $D/J=1$
  & Documento 3 & Oscillazioni non monocromatiche, leggibili \\
\bottomrule
\end{tabular}
\caption{I punti di lavoro. Il punto dinamico e il punto ``test 2'' sono
diversi, e volutamente: per la dinamica serve un segnale con più frequenze
visibili, condizione che non coincide con quella di uno stato fondamentale
interessante. Un punto unico è possibile, ma renderebbe una delle due figure
meno informativa.}
\end{table}

\FloatBarrier

\section{Come sono state ottenute le cifre di questi documenti}

Tre criteri, applicati in modo uniforme.

\begin{enumerate}[leftmargin=1.5em,itemsep=3pt]
\item \textbf{Nessun numero preso per buono dalla teoria.} Ogni valore riportato
proviene dall'esecuzione del codice. Dove una formula analitica esiste, viene
usata come controllo indipendente, non come sorgente.

\item \textbf{Due metodi indipendenti dove è possibile.} Per esempio i
correlatori sono calcolati sia con il circuito quantistico sia con
l'esponenziale di matrice diretto, che non condivide con il primo nemmeno una
riga di codice. Un accordo fra due implementazioni che condividono un bug non è
una verifica.

\item \textbf{Le discrepanze si riportano, non si smussano.} Dove il circuito e
il riferimento esatto differiscono, la differenza è mostrata e attribuita a una
causa specifica (tipicamente il numero finito di passi di Trotter), non
nascosta scegliendo una scala che la renda invisibile.
\end{enumerate}

\verifica{Come prima applicazione del criterio: lo spettro numerico ottenuto
diagonalizzando la matrice $4\times4$ coincide con la formula analitica
$E(S,M)=2JS(S{+}1)-3J+2bM$ su tutta la griglia in $b$, con scarto massimo
$8.9\times10^{-16}$ --- precisione macchina. Ogni riferimento ``esatto'' usato
nel seguito poggia su questo controllo.}

\section{Che cosa non c'è in questa serie}

Per onestà di scope, tre esclusioni esplicite:

\begin{itemize}[leftmargin=1.5em,itemsep=2pt]
\item \textbf{Il rumore.} Tutti i risultati qui sono a simulazione ideale:
vettore di stato esatto, nessun errore di gate, nessun conteggio finito. Il
sistema aperto è oggetto della Parte 2 della tesi, in sviluppo separato.
\item \textbf{L'estensione a $N=3$.} Anello e catena aperta hanno documenti
propri; qui compaiono solo come destinazione degli strumenti messi a punto.
\item \textbf{Le derivazioni.} Dove un risultato analitico serve, viene citato e
usato; la derivazione sta nei documenti di teoria già prodotti.
\end{itemize}

\end{document}
